\subsection{An attempt to derive the LVS from a 10d EFT perspective --- arXiv:2609.16117}

\textbf{Authors:} Arthur Hebecker, Andreas Schachner, Gaetano Maria Sifo

\paragraph{主な主張}
小さなサイクル上の非摂動効果が大体積極限で局所的に決まると仮定すると、10次元EFTで見積もったLVSポテンシャルの非摂動2項の体積依存性が、通常の4次元超重力の結果と一致しないことを指摘した \cite{Hebecker:2026LVSEFT}。

\paragraph{新規性}
E3 instantonとgaugino凝縮に共通する不一致を示し、幾何学的体積とカイラル座標の関係を$\tau_{s,\mathrm{hol}}=\tau_s+\alpha\log\mathcal V$で補正する場合には$\alpha=1/(2a)$が必要となり、LVSの関係式に$\mathcal O(1)$の変更が生じることを明確にした。

\paragraph{位置付け}
uplift前のLVS AdS真空を弦理論から導く際の整合性を問い直す研究であり、LVSの破綻を証明するものではなく、局所性の仮定の検証とK\"ahler座標・ポテンシャルの補正の微視的導出を課題としている。
